\documentclass[10pt,letterpaper]{article}

% - package setup
% \usepackage[doublespacing]{setspace}
\usepackage{csvsimple}
\usepackage{pgfplotstable}

\usepackage{geometry}

\usepackage{xcolor}
\definecolor{light-gray}{gray}{0.95}
\newcommand{\code}[1]{\colorbox{light-gray}{\texttt{#1}}}

% - title setup
\title{Midterm Project}
\author{Crystal Mandal}
\date{last generated \today}

\usepackage{graphics}
\begin{document}

\maketitle

\subsection*{Concept + Methodology}
The concept for this ``piece" (and, more accurately, the
messy collection of associated python and bash scripts) stems
from an interest in both the serial music concerned with pitch
classes and tone rows (e.g. Boulez and Webern) and the algorithmic
music concerned with structure and mathematical generation (e.g.
Stockhausen and Xenakis). Thus, this piece is composed of two
musical segments to be understood in counterpoint with another.
Both segments are generated the same way:

\begin{enumerate}
  \item First, a tone row and corresponding tone matrix is generated.
  \item Second, a sequence of Rows from said matrix is generated.
  \item Third, a sequence of pitches is generated using the following concatenation rule:
    \begin{center}
      Given two sequences of pitches, the second sequence is appended to the
      first at the last occurrence of the first pitch in the second sequence. Thus,
    joining [0 1 2 3 4 5] and [3 2 1] produces [0 1 2 3 2 1].
    \end{center}
  \item Finally, a sequence of durations is generated randomly exclusively for
    aural interest. The sequence of durations serves no structural purpose and is
    simply used for aural scaffolding to support the listener.
\end{enumerate}

Additionally, the metadata (that is, the sequence of rows and their transitions)
is annotated underneath the music.
The first segment, which shall be referred to as the `flat' segment, is
subject to the following constraints:

\clearpage{}

\begin{itemize}
  \item{A minimum of sixteen `measures' of music, guaranteed by a minimum of
    sixty-four generated pitches.}
  \item{There is no preference or weighting given to the sequence of rows selected;
      each possible subsequent row is equally likely, and the probability
      distribution looks `flat' - hence the name.}
  \item{The sequence of rows ends on \code{P0} or \code{I0}. }
\end{itemize}
The second segment, uses a Markov process to generate the sequence of rows and
is subject to the following constraints:

replacewith{Generated Music}
\begin{itemize}
  \item{Again, by the same methods, a minimum of sixteen `measures' and
    sixty-four pitches is generated. }
  \item{At each step in the sequence of rows, the subsequent row is chosen by a Markov
    process detailled by the Markov table in the appropriate section.}
  \item{The sequence of rows ends on \code{P0} or \code{I0}. }
\end{itemize}

This second segment is appropriately named the `Markov' segment.

\clearpage{}

\section*{Generated Music}

\subsection*{Tone Row}

The Generated Tone Row is:

\begin{quote}
{%
\parindent 0pt
\noindent
\ifx\preLilyPondExample \undefined
\else
  \expandafter\preLilyPondExample
\fi
\input{a0/lily-74006ec3-systems.tex}%
\ifx\postLilyPondExample \undefined
\else
  \expandafter\postLilyPondExample
\fi
}
\end{quote}

\subsection*{Tone Matrix}

The table generated for the tone matrix is:

% \csvautotabular[respect all]{../../python/tonematrix.csv}

\pgfplotstabletypeset[
    col sep=comma,
    string type,
]{../../python/tonematrix.csv}

\clearpage{}

\subsection*{`Flat' Segment}
The music generated with the aforementioned parameters is displayed below:

\begin{quote}
{%
\parindent 0pt
\noindent
\ifx\preLilyPondExample \undefined
\else
  \expandafter\preLilyPondExample
\fi
\input{d4/lily-415007e9-systems.tex}%
\ifx\postLilyPondExample \undefined
\else
  \expandafter\postLilyPondExample
\fi
}
\end{quote}

\newgeometry{margin = 1in}

% \vspace*{36pt}

\subsection*{Markov Chain Table}
The table generated for the Markov Chain is too long to display on one page,
so it is broken up into multiple pages below:

\vspace*{24pt}

\pgfplotstabletypeset[
    col sep=comma,
    string type,
]{../../python/markov-printfiles/seg1.csv}

\pgfplotstabletypeset[
    col sep=comma,
    string type,
]{../../python/markov-printfiles/seg2.csv}


\pgfplotstabletypeset[
    col sep=comma,
    string type,
]{../../python/markov-printfiles/seg3.csv}


\pgfplotstabletypeset[
    col sep=comma,
    string type,
]{../../python/markov-printfiles/seg4.csv}

\restoregeometry{}

\subsection*{`Markov' Segment}
The music generated by using the above Markov Table in a Markov Chain Process
is displayed below:

\begin{quote}
{%
\parindent 0pt
\noindent
\ifx\preLilyPondExample \undefined
\else
  \expandafter\preLilyPondExample
\fi
\input{d6/lily-a740e351-systems.tex}%
\ifx\postLilyPondExample \undefined
\else
  \expandafter\postLilyPondExample
\fi
}
\end{quote}

\clearpage{}

\section*{Computer Code}
The music is generated by the \code{ midterm.py } script in the
\code{ ./python } folder.

To generate a new \code{ midterm.pdf } document, run the script
\code{ ./generate.sh  }.

The generated document is available as \code{assignment-print.pdf} in the
root directory.

\end{document}
